\section{Fitting Agents to User Data from $\mathbb{X}$}
\label{sec:fitting}
In this Section, we describe the data set used to train and evaluate our approaches for imitating the types of user behavior described in the previous Section. Furthermore, the machine learning processes for fitting the models to the data is described.

\paragraph{Dataset}
\label{subsec:data}
Our experiments are based on two datasets, English and German, collected from $\mathbb{X}$. The datasets are collected around keywords concerning the political discourses in the US and Germany during the first half of 2023. The samples contain two types of content: a) Tweets (posts) from delegates of the national parliament concerning political decisions – DE: $155.000$, EN: $930.000$ and b) replies from regular $\mathbb{X}$ users towards these decisions or opinions – DE: $185.000$, EN: $17.800.000$. In addition, we labeled the topics of posts and replies using zero-shot prompt classification with Llama 3.1 70B \cite{zheng2023judging}. This yields two distinct agent types, the a) political actor for posts and the b) unspecified $\mathbb{X}$ user for replies, both across two languages. 

\paragraph{Preprocessing}
\label{subsec:preprocessing}
We apply moderate content filtering to enhance the quality of the selected data:

\emph{External Source:} We remove samples that include URLs to external sources. Although this significantly reduces the number of samples, our models cannot access this information during the training and inference process.

\emph{Retweets:} Tweets starting with \textit{RT} mark a shared post. The author does not add new content, but replicates the original text and shares it with their followers. We exclude these samples because we cannot verify if the shared content is in line with the retweeting author's opinion. 

\emph{Content Length:} We aim to focus on content that contains arguments or opinions, particularly given the later applied discourse metrics. Thus, we remove samples with less than $32$ characters based on our heuristic assumption that a specified length is necessary to convey an opinion.

After the initial pre-processing, we further reduce both datasets to retrieve only the most active users measured by the number of posts and replies. This is required for modeling both tasks on a user-based level by providing user examples (few-shot) during prompting.

\paragraph{Content Generation Alignment Pipeline}
\label{subsec:training}
We train two model adapters on top of Llama-3.2-3B-Instruct \citep{dubey2024llama} using the supervised fine-tuning (SFT) paradigm through the transformer reinforcement learning (TRL) library using the default pipeline. The trained adapters are unmodified LoRA \citep{hu2021lora} matrices provided by the parameter-efficient fine-tuning (PEFT) package. We optimize for two objectives (posting and replying) with contents from the original users as the labeled data.

\paragraph{Behavior Likelihood Alignment Pipeline}
For the decision task, we train our dual input model using embeddings generated by a frozen BERT-based encoder that was pretrained to capture semantic representations of tweets \citep{zhang2023twhin}. We optimize with AdamW \citep{loshchilov2017fixing} using the same data set as for the generative tasks, making the assumption that when a user has not commented on tweets in our collection, this represents an active decision not to reply. To maintain balance, we sample an equal number of negative and positive instances for both languages.

\paragraph{Data and Code Availability}
The complete technical pipeline for this study is publicly available on GitHub at \url{https://anonymous.4open.science/r/TWON-Agents-D0E2}. This repository includes the source code for all the data processing, training, and evaluation steps described in this paper. The raw and processed datasets used in this study will be made available after acceptance.

% (anonymized during review) at \url{https://anonymous.4open.science/r/TWON-Agents-D0E2/}